\documentclass[11pt]{article}
\usepackage[a4paper,margin=27mm]{geometry}
\usepackage{amsmath,amssymb,amsthm,booktabs,hyperref}
\newtheorem{theorem}{Theorem}
\newtheorem{lemma}{Lemma}
\newcommand{\R}{\mathbb R}
\title{The shared square register cannot be removed by Laurent reweighting}
\author{\(\mathbb Z\mathbb Q\) Diophantine Project}
\date{23 November 2026}
\begin{document}\maketitle
\begin{abstract}
For each of the retained 163- and 179-coordinate quartics, deleting the
single register for the square of the scalar normalizer cannot be achieved
by any diagonal Laurent reweighting in the two existing units while the
physical equations and other unit registers are retained. The obstruction
already holds for real weights. Six source rows give a short convex
inequality proof; independently checkable integer certificates give the
same conclusion. This is a restriction on the specified compilation,
not a lower bound for quartics defining the integers.
\end{abstract}

\paragraph{Attribution and scope.}
The shared-product constructions originate in the notes of Gabe Conant
with Astra 6.2. The scalar-point construction and current physical source
rows are those of Jose Brox's project continuation v82. The present
exclusion is a subsequent deduction. It does not alter any arithmetic
premise or polynomial record. Its source is
\path{papers/quartic/scalar_point163_179_20261120_v82.tex}
and the corresponding \texttt{source163.json} and
\texttt{source179.json} physical tables.

\section{The compilation model}
Let \(c,R\ne0\) be the two units. Remove the register \(R_2=R^2\)
and its defining row. Keep all remaining physical variables, equations,
and the unit registers for \(c,c^{-1},c^{-2},R\); the 179-coordinate
source also keeps \(R^{-1}\). A diagonal Laurent reweighting replaces
each remaining physical variable \(v_i\) by
\(c^{a_i}R^{b_i}y_i\) and divides each row by a fixed monomial
\(c^{A_j}R^{B_j}\). It must represent every resulting term by a
polynomial of total degree at most two using the retained registers.
No row combinations, nonlinear changes, or additional registers are
allowed. Fixed nonzero rational constants and powers of the fixed
coefficient parameter have weight zero. Actual Laurent exponents are
integral; allowing real weight vectors only relaxes the necessary conditions.

Put \(e_c=(1,0)\), \(e_R=(0,1)\). The convex hull of exponent
vectors attainable with one unit slot is
\[
 P_{163}=\operatorname{conv}\{(-2,0),(1,0),(0,1)\},\qquad
 P_{179}=\operatorname{conv}\{(-2,0),(1,0),(0,1),(0,-1)\}.
\]
The exponent \((0,0)\) corresponds to an unused slot. If a term
already contains two physical coordinates, its residual unit exponent
must be zero. A term with no physical coordinate must have residual
exponent in \(2P_m\). These are necessary conditions, whether or not
the weights are integral.

\section{Six physical rows}
Use the source-coordinate names
\[
 F=v_{65},\quad J_0=v_{109},\quad J_1=v_{110},\quad
 T=v_{111},\quad \sigma=v_{104},\quad B=v_{105},\quad
 E=v_{11},\quad P=v_{193}.
\]
The two sources contain the following terms with nonzero coefficients.
Only the indicated terms of each row are needed.
\begin{center}\small
\begin{tabular}{ll}
\toprule
Source row & Selected terms, up to nonzero rational constants\\
\midrule
\texttt{new.det0.factor0.0}&\(1,\ c^2F^2\)\\
\texttt{v82.point.product.0}&\(c^{-1}FP,\ J_1P\)\\
\texttt{v82.scalar.partial\_trace}&\(J_0^2,\ J_1^2,\ J_0T\)\\
\texttt{v82.scalar.middle}&\(E\sigma,\ EB\)\\
\texttt{v82.scalar.input}&\(B\sigma,\ R\)\\
\texttt{v82.scalar.inverse}&\(T\sigma,\ R^4\)\\
\bottomrule
\end{tabular}
\end{center}
In the last two rows the eliminated \(R_2\) has been replaced literally
by \(R^2\). All selected physical monomials remain distinct, so none
can cancel under a diagonal rescaling.

\begin{lemma}
If such a reweighting exists, there are vectors \(t,s\in\R^2\)
satisfying
\begin{equation}\label{ineq}
 -t-2e_c\in P_m,\qquad
 e_R-2s\in2P_m,\qquad
 4e_R-t-s\in2P_m.
\end{equation}
\end{lemma}
\begin{proof}
Write \(w_i\) for the weight of a physical coordinate and \(r_j\)
for that of a row. Quadratic terms must have residual exponent zero.
The point-product row therefore gives \(w_{J_1}=w_F-e_c\).
The partial-trace row gives
\(w_{J_0}=w_{J_1}=w_T=:t\), so \(w_F=t+e_c\).
The detector row has weight \(2w_F+2e_c\). Its constant term
therefore requires \(-2w_F-2e_c\in2P_m\), or
\(-t-2e_c\in P_m\).

The middle pole row gives \(w_\sigma=w_B=:s\). The input
row has weight \(2s\); its linear \(R\) term requires
\(e_R-2s\in2P_m\). The inverse row has weight \(t+s\);
its \(R^4\) term requires \(4e_R-t-s\in2P_m\).
This proves all three necessary conditions.
\end{proof}

\begin{theorem}
Neither retained general quartic admits a termwise quadratic compilation
of the stated kind after deleting its \(R^2\) register. This remains
true when both unit weights and all row weights are allowed to vary.
\end{theorem}
\begin{proof}
For the 163-coordinate source, the second coordinate of every point of
\(P_{163}\) lies in \([0,1]\). Apply this linear functional to
\eqref{ineq}: it gives \(t_R\le0\), \(s_R\le1/2\), and
\(t_R+s_R\ge2\), a contradiction.

For the 179-coordinate source, use the linear functional
\(\ell(a,b)=a+2b\). Its minimum and maximum on \(P_{179}\)
are \(-2\) and \(2\). The first relation in \eqref{ineq} gives
\(-\ell(t)-2\ge-2\), hence \(\ell(t)\le0\).
The second gives \(2-2\ell(s)\ge-4\), hence \(\ell(s)\le3\).
The third gives \(8-\ell(t)-\ell(s)\le4\), hence
\(\ell(t)+\ell(s)\ge4\), again a contradiction.
\end{proof}

\section{Exact certificate and consequences}
The accompanying standard-library checker reconstructs the available
unit-exponent polygons, binds each selected inequality to a nonzero term
of the physical source, and sums integer inequalities exactly. The
19-term certificates give respectively \(0\le-9\) and \(0\le-6\).
No optimization program is required to verify them. The search that
discovered the certificates used linear programming, but the proof and
final verification use only exact integer arithmetic.

The obstruction identifies a specific limitation of reweighting. A
smaller quartic might still follow from nonlinear coordinate changes,
variable-dependent combinations of equations, a different native pole
presentation, different point charts, or a different arithmetic source.
None of these possibilities is excluded by this theorem. In particular
the theorem does not prove a global minimum of 163 or 179 variables.
\end{document}
